\section{Опис реалізації}

\subsection{Підготовка корпусу}

Скрипт \code{tools/prepare\_dataset.py} читає книжки з архіву у відсортованому порядку,
вирізає ліцензійні шапку й футер Project Gutenberg (між маркерами
\code{*** START OF ...} і \code{*** END OF ...}), нормалізує кінці рядків і пакує абзаци в
документи розміром близько 100~КБ. Абзац, довший за ліміт, ріжеться за останнім переведенням
рядка або пробілом до межі, тож слова не розриваються; короткий «хвіст» книжки (менше чверті
ліміту) приєднується до попереднього документа. Скрипт детермінований: повторний запуск дає
ті самі файли. У Docker Compose його виконує одноразовий сервіс \code{dataset-init}, записуючи
результат в іменований том \code{corpus} на файловій системі Linux.

Розміщення корпусу в томі, а не в каталозі диска Windows, -- результат першого ж вимірювання.
Коли сервер читав документи безпосередньо з \code{D:\textbackslash} через механізм спільних
каталогів Docker Desktop (міст Windows $\leftrightarrow$ WSL2), побудова на 12 потоках
тривала 42,3~с, з яких 19,3~с займав лише обхід каталогу, а сумарний час потоків на читання
файлів становив 238~с: обчислення фактично стояли в черзі до файлової системи. Після
перенесення корпусу в том обхід каталогу скоротився до 0,13~с, а вся побудова -- до 4,0~с.
Характеристики отриманого корпусу наведено в табл.~\ref{tab:corpus}.

\begin{table}[H]
\caption{Характеристики корпусу}\label{tab:corpus}
\centering\begin{tabularx}{\textwidth}{Xr}
\toprule
Показник & Значення\\
\midrule
Книжок в архіві D1.7GB.zip & 2\,809\\
Документів після розрізання & 16\,793\\
Обсяг документів & 1\,623\,765\,577 байт (1,62~ГБ)\\
Цільовий розмір документа & $\approx$100~КБ\\
Унікальних термів у словнику & 840\,954\\
Записів у списках входжень (пар «терм -- документ») & 33\,197\,116\\
Середня довжина документа після токенізації & 5\,501 терм\\
\bottomrule
\end{tabularx}
\end{table}

\subsection{Пул потоків}

Клас \code{ThreadPool} -- фіксований набір робочих потоків і необмежена FIFO-черга задач
\code{ArrayDeque<Runnable>}, захищена одним \code{ReentrantLock}. Вільні потоки чекають на
умові \code{notEmpty} у циклі з перевіркою предиката, що захищає від хибних пробуджень і
виключає активне очікування. Метод \code{submit} загортає задачу в \code{FutureTask} і
повертає \code{Future}: виняток задачі зберігається в ньому й повертається викликачу через
\code{ExecutionException}, а робочий потік продовжує роботу. Стан пулу змінюється за схемою
\code{RUNNING} $\rightarrow$ \code{SHUTDOWN} $\rightarrow$ \code{TERMINATED}: \code{shutdown()}
перестає приймати задачі, але дає виконати вже поставлені; \code{shutdownNow()} додатково
повертає невиконані задачі й перериває потоки; \code{close()} виконує коректне завершення з
приєднанням усіх потоків, що дозволяє використовувати пул у конструкції try-with-resources.

Інваріанти пулу: черга і стан читаються та змінюються лише під замком; задача виймається з
черги рівно одним потоком, тож виконується рівно один раз; \code{Future.get()} встановлює
відношення happens-before між діями задачі та кодом, який дочекався її результату [15]. Ця
остання властивість використовується в побудові індексу для безпечної передачі даних між
фазами без додаткових бар'єрів. У сервері працюють три види пулів: пул HTTP-з'єднань (64
потоки), пул WebSocket-сеансів (32 потоки) і пул побудови індексу, який створюється на кожен
запит індексації з кількістю потоків \code{threads}.

\subsection{Структура інвертованого індексу}

Інвертований індекс зіставляє кожному терму список документів, у яких він трапляється [3].
У реалізації кожен запис списку (posting) упаковано в одне 64-бітне число: старші 32 біти --
ідентифікатор документа, молодші -- частота терма в ньому. Упаковка вдвічі зменшує кількість
об'єктів порівняно з парою масивів і перетворює сортування за документом на звичайне
\code{Arrays.sort(long[])}, оскільки ідентифікатори невід'ємні. Словник поділено на $P$
хеш-партицій (за замовчуванням 64); партиція терма визначається хешуванням Фібоначчі
$h \cdot \mathtt{0x9E3779B9}$ з подальшим взяттям залишку, що рівномірно розподіляє схожі
рядки. Кожна партиція -- звичайна \code{HashMap<String, long[]>}.

Готовий \code{InvertedIndex} незмінний: після побудови жодна партиція не модифікується. Він
публікується разом зі статистикою побудови як об'єкт \code{Snapshot} через
\code{AtomicReference}. Запис у volatile-посилання після завершення всіх фаз і його читання
пошуковим потоком утворюють пару release/acquire, тому всі записи, зроблені під час побудови,
видимі читачеві без жодних блокувань [15]. Пошуковий запит один раз читає посилання і далі
працює з цим знімком до кінця, навіть якщо тим часом паралельна перебудова опублікувала
новий (схема build-then-swap). Ціною є те, що під час перебудови в пам'яті тимчасово
існують два індекси (підрозділ~3.10).

\subsection{Паралельна побудова індексу}

Побудова складається з п'яти фаз (рис.~\ref{fig:seq_index.png}):
\begin{enumerate}[label=\arabic*)]
\item \textit{list} (викликаючий потік) -- рекурсивний обхід каталогу, відбір \code{*.txt},
сортування шляхів; позиція у списку є ідентифікатором документа;
\item \textit{map} (пул) -- список ріжеться на батчі по 32 документи, тобто 525 задач для
16\,793 документів; задача читає файли і токенізує їх у приватну таблицю термів, нічого не
синхронізуючи;
\item \textit{merge} (пул, та сама задача) -- приватна таблиця передається будівнику
\code{IndexBuilder.addBatch}, який зливає її у спільну структуру;
\item \textit{finalize} (пул) -- по одній задачі на партицію сортує списки входжень і
перетворює їх на компактні масиви;
\item \textit{swap} -- новий знімок атомарно замінює попередній.
\end{enumerate}

Фіксований розмір батча зроблено свідомо: за будь-якої кількості потоків виконується однаковий
набір задач, тож у вимірюваннях змінюється лише паралелізм, а не гранулярність. 525 задач на
16 потоків -- понад 30 задач на потік, чого достатньо для балансування навантаження, коли
документи мають різну кількість унікальних слів. Довжини документів записуються в масив
\code{DocumentTable} без синхронізації: кожну комірку пише рівно одна задача, а
\code{Future.get()} робить записи видимими потоку, що виконує фінальну фазу. Неможливість
прочитати файл не перериває побудову: документ отримує нульову довжину і не враховується в
\code{documents\_indexed}; якщо не прочитано жодного документа, побудова відхиляється, а
чинний індекс зберігається.

\subsection{Неблокуюче злиття: партиції з CAS-скриньками}

\code{PartitionedIndexBuilder} -- основна реалізація спільної структури. Кожна з $P$ партицій
має приватну \code{HashMap} термів, неблокуючу «скриньку» (стек Трайбера [7]) і прапорець
власника \code{busy} типу \code{AtomicBoolean}. Задача, що завершила свій батч, розбиває його
на фрагменти за партиціями і для кожного фрагмента виконує такі кроки:
\begin{enumerate}[label=\arabic*)]
\item кладе фрагмент у скриньку партиції (операція CAS на вершині стека);
\item намагається стати власником партиції: \code{busy.compareAndSet(false, true)};
\item якщо вдалося -- вибирає зі скриньки всі фрагменти (свій і, можливо, чужі), зливає їх
у \code{HashMap} партиції, знімає прапорець і \textit{повторно перевіряє скриньку};
\item якщо партиція зайнята -- не чекає, а переходить до наступного фрагмента: її фрагмент
зіллє поточний власник (передача, hand-off).
\end{enumerate}

Така схема є різновидом «комбінування» (flat combining [8]): потік, що отримав доступ до
структури, виконує операції і за інших. Жоден потік ніколи не блокується на партиції.
Повторна перевірка після зняття прапорця закриває вікно, у якому фрагмент потрапив у скриньку
вже після останнього \code{pop} власника, але до \code{busy.set(false)}: той, хто поклав
фрагмент, бачить партицію зайнятою і йде далі, а власник після звільнення помічає непорожню
скриньку і знову бере її. Оскільки кожен, хто кладе фрагмент, після цього обов'язково
намагається стати власником, після повернення всіх викликів \code{addBatch} скриньки
порожні; фінальна фаза для певності виконує ще одне спорожнення.

Відсутність стану гонки обґрунтовується моделлю пам'яті Java [15]. \code{HashMap} партиції
змінює лише поточний власник. Право власності передається через volatile-прапорець:
\code{set(false)} має семантику release, успішний \code{compareAndSet} -- acquire, тому
кожен наступний власник бачить усі зміни попереднього. Фрагменти потрапляють від автора до
власника через CAS на вершині стека, що також встановлює happens-before. Проблема ABA для
стека Трайбера в Java не виникає: вузол не може бути повторно використаний, доки на нього
посилається хоч один потік, бо пам'ять звільняє збирач сміття. Взаємне блокування
неможливе, бо замків немає взагалі.

Важливо, що неблокуючий протокол прибирає \textit{очікування}, але не \textit{серіалізацію}:
злиття в одну партицію все одно виконується послідовно, тільки тим потоком, який першим
захопив власність. Наслідки цього видно з експерименту з однією партицією (підрозділ~3.7).

\subsection{Блокувальна альтернатива: lock striping}

Для порівняння реалізовано \code{StripedIndexBuilder} -- класичну схему розщеплення замків:
словник поділено на $S$ смуг, кожна з яких -- \code{HashMap} під власним \code{ReentrantLock}.
Задача групує батч за смугами і зливає кожну групу, утримуючи замок відповідної смуги. Потік
у кожен момент утримує щонайбільше один замок, тому кругове очікування, а з ним і взаємне
блокування, неможливі. За $S = 1$ схема вироджується в один глобальний замок. Будівник
вимірює кількість захоплень, кількість випадків, коли замок був зайнятий (спершу
\code{tryLock}, потім блокувальний \code{lock}), і сумарний час очікування -- ці лічильники
використано в підрозділі~3.7.

\subsection{Токенізація: дві версії}

Токен -- це максимальна послідовність ASCII-літер і цифр довжиною від 2 до 32 символів,
приведена до нижнього регістру; довші послідовності (хеші, злиті рядки) відкидаються, будь-який
інший байт, зокрема не-ASCII, є роздільником. Англійські службові слова (132 слова) у
індекс не потрапляють. Ті самі правила застосовуються до тексту запиту і до термів від моделі.

\textit{Перша версія} токенізатора створювала рядок \code{String} для кожного слова і
підраховувала частоти в \code{HashMap}, що очищалася для кожного документа. Для 1,6~ГБ
тексту це сотні мільйонів короткоживучих об'єктів. Розклад часу за фазами (підрозділ~3.5)
показав, що сумарний час токенізації зростає з кількістю потоків значно швидше, ніж мав би:
потоки конкурують за спільні ресурси -- виділення пам'яті, роботу збирача сміття, шину пам'яті.

\textit{Друга версія} (\code{BatchTermTable}) хешує байти слова на місці, не копіюючи їх, і
шукає терм у таблиці з відкритою адресацією та лінійним зондуванням, що належить задачі.
Хеш обчислюється тією самою формулою $h = 31h + c$, що й \code{String.hashCode()}, тож
порівняння з уже збереженим рядком виконується без перетворень. Рядок створюється лише тоді,
коли терм уперше трапляється в батчі. Частоти накопичуються в парі полів
«поточний документ / лічильник» запису і скидаються в список входжень при зміні документа,
тому окремої таблиці на документ не існує. Стару версію залишено як параметр
\code{"mapper": "hashmap"}; тести перевіряють, що обидві версії будують ідентичний індекс.
Ефект заміни наведено в підрозділі~3.6.

\subsection{HTTP-сервер і паралельна обробка запитів}

\code{HttpServer} написано поверх блокувальних сокетів. Окремий потік виконує цикл
\code{accept} і передає кожне з'єднання в пул HTTP-потоків. Потік пулу обслуговує з'єднання
протягом усього його життя: читає запит, викликає обробник, записує відповідь і, якщо клієнт
не просив закрити з'єднання, читає наступний запит (HTTP/1.1 keep-alive [14]). Повторне
використання TCP-з'єднання прибирає з кожного запиту рукостискання і є головною причиною
високої пропускної здатності в підрозділі~3.8. Тайм-аут простою з'єднання -- 15~с, після
10\,000 запитів сервер закриває з'єднання заголовком \code{Connection: close}, щоб один клієнт
не утримував потік безстроково. Запит на WebSocket-рукостискання переводить сокет у пул
WebSocket-сеансів, тож тривалі сеанси не займають HTTP-потоки.

Сервер обмежує вхідні дані: рядок запиту із заголовками -- до 16~КіБ (431), тіло -- до 1~МіБ
(413), тіло з \code{Transfer-Encoding} не підтримується (501). Невідомий шлях дає 404,
невідповідний метод -- 405 із заголовком \code{Allow}, звичайний GET на \code{/ws/search} --
426. Невідомі поля JSON ігноруються із записом у журнал. Під час зупинки (\code{docker stop}
надсилає SIGTERM, який обробляє shutdown hook JVM) сервер закриває слухаючий сокет і
простоюючі keep-alive з'єднання, а запитам, що вже виконуються, дає до 8~с на завершення.

\subsection{WebSocket}

Протокол RFC~6455 [13] реалізовано без сторонніх бібліотек. Рукостискання перевіряє метод
GET, заголовки \code{Connection: Upgrade}, \code{Sec-WebSocket-Key} (16 байт у Base64) і версію
13; відповідь містить \code{Sec-WebSocket-Accept = Base64(SHA-1(key + GUID))}, що перевірено на
прикладі з самого RFC. Читач кадрів підтримує довжини 7/16/64 біти, демаскування, ping/pong і
\textit{фрагментовані повідомлення} (кадри продовження). Порушення протоколу закривають
з'єднання з відповідним кодом: зарезервовані біти, незамаскований кадр, керуючий кадр понад
125 байтів, продовження без початого повідомлення -- 1002; бінарне повідомлення -- 1003;
некоректний UTF-8 -- 1007; повідомлення понад 64~КіБ -- 1009. Запис кадрів серіалізовано
замком, щоб pong, надісланий під час запису великого повідомлення, не вклинився в нього.
Закриття виконується за повним протоколом: сервер надсилає close-кадр і до 1~с чекає на
зустрічний.

\subsection{Ранжування та сніпети}

Документи ранжуються функцією Okapi BM25 [9]:
\begin{equation}
\mathrm{score}(d) = \sum_{t \in q} w_t \cdot \mathrm{idf}(t) \cdot
\frac{tf_{t,d}\,(k_1 + 1)}{tf_{t,d} + k_1\left(1 - b + b\,\dfrac{|d|}{\mathrm{avgdl}}\right)},
\end{equation}
\begin{equation}
\mathrm{idf}(t) = \ln\!\left(1 + \frac{N - df_t + 0{,}5}{df_t + 0{,}5}\right),
\end{equation}
де $k_1 = 1{,}2$, $b = 0{,}75$, $N$ -- кількість документів, $df_t$ -- документна частота
терма, $|d|$ -- довжина документа в термах. Вага $w_t$ дорівнює 1 для слів користувача і 0,4 для
слів, доданих моделлю: розширення підвищує повноту, але не може переважити те, що
користувач написав сам. На відміну від простого підрахунку збігів, BM25 враховує
рідкісність терма і частоту його вживання з поправкою на довжину документа, тож документ,
де «whale» трапляється сто разів, випереджає документ із випадковою згадкою.

Обчислення не потребує синхронізації: кожен запит виділяє власний масив-акумулятор на $N$
елементів, проходить списки входжень своїх термів і вибирає 10 найкращих документів
мін-купою. Рівні оцінки впорядковуються за ідентифікатором документа, тому видача
детермінована. Поле \code{score} у відповіді -- оцінка, поділена на найкращу оцінку запиту
(діапазон $(0; 1]$), сире значення BM25 повертається в полі \code{bm25}.

Сніпет будується читанням файла блоками по 64~КіБ до першого входження найрідкіснішого
знайденого слова користувача (якщо таких немає -- слова розширення). Збіг зараховується лише
за цілим словом, тому \code{parallel} не знаходиться всередині \code{unparalleled}, а сніпет
узгоджений з ранжуванням. Вікно до 180 символів зсувається до меж слів, пробіли згортаються,
не-ASCII байти замінюються символом-заповнювачем.

\subsection{Розширення запитів через LLM}

\code{QueryExpander} звертається до Ollama \code{/api/generate} без стрімінгу, з форматом
відповіді JSON, температурою 0,2 і лімітом 128 токенів. Промпт вимагає 5--8 англійських
синонімів або близьких ключових слів окремими словами, без повторів і без слів самого
запиту, і містить приклад «запит -- відповідь» на сторонньому запиті. Оскільки модель на CPU
є найповільнішою ланкою конвеєра, навколо виклику побудовано чотири механізми:
\begin{itemize}
\item \textit{LRU-кеш} останніх 1\,024 розширень (\code{LinkedHashMap} у порядку доступу під
замком);
\item \textit{об'єднання однакових запитів} (single-flight): одночасні запити з тим самим
набором термів чекають на один спільний \code{CompletableFuture}, зареєстрований атомарним
\code{putIfAbsent} у \code{ConcurrentHashMap};
\item \textit{обмеження паралельності}: не більше двох одночасних викликів моделі
(\code{Semaphore}, узгоджено з \code{OLLAMA\_NUM\_PARALLEL=2});
\item \textit{дедлайн}: на очікування слота, сам виклик і очікування спільного результату
разом дається 20~с; після тайм-ауту чи будь-якої помилки моделі пошук виконується за
початковим запитом, а у відповіді вказується джерело \code{timeout} або \code{error}.
\end{itemize}

Відповідь моделі розбирається захищено: прибираються обгортки коду, вирізається JSON-об'єкт
або масив з будь-якого оточення, а якщо відповідь обірвана лімітом токенів, терми
відновлюються регулярним виразом за закритими лапками. Отримані слова проходять той самий
токенізатор, що й індекс, з них прибираються дублікати, службові слова, числа і слова запиту;
залишається не більше 8 термів.

\subsection{Клієнти}

\textit{CLI-клієнт} (\code{SearchCli}) працює всередині контейнера, тому на хості не
потрібна Java: обгортки \code{./search} (Git Bash, Linux) і \code{search.cmd} (Windows)
викликають \code{docker compose exec server search}. Клієнт підтримує пошук з розширенням і
без, кількість результатів, побудову індексу та вивід статистики; у терміналі знайдені слова
підсвічуються кольором.

\textit{Web UI} -- одна сторінка на HTML, CSS і JavaScript без фреймворків, яку роздає сам
сервер. Пошук іде через WebSocket: три картки етапів конвеєра підсвічуються в міру надходження
повідомлень, показують час кожного етапу і джерело розширення, слова користувача та моделі
відображаються різними кольорами, у сніпетах вони виділяються. З'єднання перевідкривається
після розриву. Окрема панель будує індекс із заданою кількістю потоків і реалізацією та
показує статистику. Усі дані сервера вставляються в сторінку лише через \code{textContent} і
вузли DOM, без \code{innerHTML}, що виключає XSS через вміст документів.

На рис.~\ref{fig:webui.png} показано пошук «sherlock holmes» з розширенням: модель додала
терми «consulting», «detective», «baker», «street» (позначені синім), розширення тривало
734~мс, пошук за індексом -- 9~мс. Той самий запит без розширення (рис.~\ref{fig:webui_plain.png})
виконується за 2~мс; порядок результатів змінюється, бо документи, де трапляються також
слова розширення, отримують додаткову оцінку.

\fig{webui.png}{Web UI: пошук з розширенням через LLM}
\fig{webui_plain.png}{Web UI: той самий запит без розширення}
\IfFileExists{figures/cli.png}{\fig{cli.png}{CLI-клієнт: пошук з розширенням}}{}

\subsection{Контейнеризація}

\code{Dockerfile} багатоетапний: на етапі збирання в образі \code{eclipse-temurin:21-jdk}
скрипт \code{build.sh} компілює код, пакує jar і запускає всі тести -- якщо хоч один тест не
проходить, образ не збирається. Робочий образ на базі \code{eclipse-temurin:21-jre} містить
лише jar, статичні файли і скрипти. JVM запускається з
\code{-XX:MaxRAMPercentage=75 -XX:+UseG1GC} і стає процесом PID~1, тому отримує SIGTERM від
\code{docker stop}. Перевірка стану контейнера виконується без curl -- запитом через
\code{/dev/tcp} у bash.
